% sec10_3.tex — Section 10.3: Testing the Null of No Cointegration

Having introduced the notion of cointegration, we need to deal with two issues.
The first is how to determine the cointegration rank, and the second is how to
estimate and do inference on cointegrating vectors. The first will be discussed in
this section, and the second will be the topic of the next section. There are several
procedures for determining the cointegration rank. Among them are Johansen's
(1988) likelihood ratio test derived from the maximum likelihood estimation of
the VECM (briefly covered at the end of previous section) and the common trend
procedure of Stock and Watson (1988). These procedures allow us to test the null
of $h = h_0$ where $h_0$ is some arbitrary integer between 0 and $n - 1$. We will not
cover these procedures.\footnote{%
See Maddala and Kim (1998, Section~7) for a catalogue of available procedures.}
Here, we cover only the simple test suggested by Engle
and Granger (1987) and extended by Phillips and Ouliaris (1990). In that test, the
null hypothesis is that $h = 0$ (no cointegration) and the alternative is that $h \geq 1$.

\subsection*{Spurious Regressions}

The test of Engle and Granger (1987) is based on OLS estimation of the regression
\begin{equation}
  y_{1t} = \mu + \boldsymbol{\gamma}' \mathbf{y}_{2t} + z_t^*,
  \label{eq:10.3.1} \tag{10.3.1}
\end{equation}
where $y_{1t}$ is the first element of $\mathbf{y}_t$, $\mathbf{y}_{2t}$ is the vector of the remaining $n - 1$ ele\-ments,
and $z_t^*$ is an error term. This regression would be a cointegrating regression
if $h = 1$ and $y_{1t}$ were part of the cointegrating relationship. Under the null of $h = 0$
(no cointegration), however, this regression does not represent a cointegrating rela\-tionship.
Let $(\hat{\mu}, \hat{\boldsymbol{\gamma}})$ be the OLS coefficient estimates of $(\mu, \boldsymbol{\gamma})$. It turns out that $\hat{\boldsymbol{\gamma}}$
does not provide consistent estimates of any population parameters of the system!
For example, even if $y_{1t}$ is unrelated to $\mathbf{y}_{2t}$ (in that $\Delta y_{1t}$ and $\Delta \mathbf{y}_{2s}$ are independ\-ent
for all $s, t$), the $t$- and $F$-statistics associated with the OLS estimates become
arbitrarily large as the sample size increases, giving a false impression that there
is a close connection between $y_{1t}$ and $\mathbf{y}_{2t}$. This phenomenon, called the \textbf{spurious
regression}, was first discovered in Monte Carlo experiments by Granger and New\-bold
(1974). Phillips (1986) theoretically derived the large-sample distributions of
the statistics for spurious regressions. For example, the $t$-value, if divided by $\sqrt{T}$,
converges to a nondegenerate distribution.

\subsection*{The Residual-Based Test for Cointegration}

The regression (10.3.1) nevertheless provides a useful device for testing the null
of no cointegration, because the OLS residuals, $y_{1t} - \hat{\mu} - \hat{\boldsymbol{\gamma}}'\mathbf{y}_{2t}$, should appear to
have a stochastic trend if $\mathbf{y}_t$ is not cointegrated and be stationary otherwise. Engle
and Granger (1987) suggested applying the ADF $t$-test to the residuals in order to
test the null of no cointegration. Because of the use of the residuals, the test is
called the \textbf{residual-based test for cointegration}. The asymptotic distributions of
the test statistic for some leading unit-root tests were derived theoretically and the
(asymptotic) critical values tabulated by Phillips and Ouliaris (1990) and Hansen
(1992a).

In contrast to the univariate unit-root tests, the asymptotic distributions (and
hence the asymptotic critical values) depend on the dimension $n$ of the system. This
is because the residuals depend on $(\hat{\mu}, \hat{\boldsymbol{\gamma}})$, which, being estimates based on data,
are random variables. Here we indicate the appropriate asymptotic critical values
when the unit-root test applied to the residuals is the ADF $t$-test of Proposition~9.6
for autoregressions without a constant or time. That is, the ADF $t$-statistic is the $t$-value
on the $x_{t-1}$ coefficient in the following augmented autoregression estimated
on residuals:
\begin{equation}
  \Delta x_t = (\rho - 1) x_{t-1}
    + \zeta_1 \Delta x_{t-1} + \cdots + \zeta_p \Delta x_{t-p}
    + \text{error},
  \label{eq:10.3.2} \tag{10.3.2}
\end{equation}
where $x_t$ here is the residual from regression (10.3.1). There is no need to include
a constant in this augmented autoregression because, with the regression already
including a constant, the sample mean of the residuals is guaranteed to be zero.
There is no need to include time either, because the variables of the regression
(10.3.1) implicitly or explicitly include time trends (see below for more on this).
The number of lagged changes, $p$, needs to be increased to infinity with the sample
size $T$, but at a rate slower than $T^{1/3}$. More precisely,\footnote{%
See Phillips and Ouliaris (1990, Theorem~4.2). The lag length $p$ can be a random variable because it can
be data-dependent. This condition is the same as in the Said-Dickey extension of the ADF tests in the previous
chapter (see Section~9.4), except that $p$ here can be a random variable.}
\begin{equation}
  p \to \infty \quad \text{but} \quad
  \frac{p}{T^{1/3}} \pto 0
  \quad \text{as } T \to \infty.
  \label{eq:10.3.3} \tag{10.3.3}
\end{equation}
The critical values are the same if Phillips' $Z_t$-test (see Analytical Exercise~6 of the
previous chapter) is used instead of the ADF~$t$. There are three cases to consider.

\begin{enumerate}
\item $\E(\Delta \mathbf{y}_{2t}) = \mathbf{0}$ and $\E(\Delta y_{1t}) = 0$, so no elements of the I(1) system have drift.
  The appropriate critical values are in Table~10.1(a), which reproduces Table
  IIb of Phillips and Ouliaris (1990). For a statement of the conditions on the
  VMA representation under which the asymptotic distribution is derived, see
  Hamilton (1994, Proposition~19.4).

\item $\E(\Delta \mathbf{y}_{2t}) \neq \mathbf{0}$ but $\E(\Delta y_{1t})$ may or may not be zero. This case was discussed by
  Hansen (1992a). Let $g$ ($\equiv n - 1$) be the number of regressors besides a con\-stant
  in the regression (10.3.1). Some of the $g$ regressors have drift. Since lin\-ear
  trends from different regressors can be combined into one,\footnote{%
  For example, let $n = 3$ so that there are two regressors, $y_{2t}$ and $y_{3t}$ with coefficients $\gamma_1$ and $\gamma_2$, respectively.
  If $\delta_2$ and $\delta_3$ are the drifts in $y_{2t}$ and $y_{3t}$, respectively, then the linear trends in regression (10.3.1) are $\gamma_1\delta_2 t$ and
  $\gamma_2\delta_3 t$, which can be combined into a single time trend $(\gamma_1\delta_2 + \gamma_2\delta_3)t$.}
  the regression
  (10.3.1) can be rewritten as a regression of $y_{1t}$ on a constant, $g - 1$ I(1) regres\-sors
  without drift, and one I(1) regressor with drift. Now, since linear trends
  dominate stochastic trends (in the sense made precise in the discussion of the
  BN decomposition in the previous chapter), the I(1) regressor with a trend
  behaves very much like time in large samples. So the residuals are ``asymp\-totically
  the same'' as the residuals from a regression with a constant, $g - 1$
  driftless I(1) variables, and time as regressors, in the sense that the limiting
  distribution of a statistic based on the residuals from the former regression is
  the same as that from the latter regression.

  The critical values for the ADF $t$ test based on the latter regression are
  tabulated in Table~IIc of Phillips and Ouliaris (1990). Therefore, to find the
  appropriate critical value when the regression (10.3.1) has a constant and $g$
  regressors but not time, turn to this table for $g - 1$ regressors. If the regression
  has only one regressor besides the constant (i.e., if $g = 1$), then the regres\-sion
  is asymptotically equivalent to a regression of $y_{1t}$ on a constant and time.
  For this case, the limiting distribution of the ADF $t$-statistic calculated from
  the residuals turns out to be the Dickey-Fuller distribution ($DF_t^\tau$) of Proposi\-tion~9.8.
  Table~10.1(b) combines these distributions: for the case of one
  regressor ($g = 1$), it shows the ADF $t^\tau$-distribution, and for the case of $g$ ($> 1$)
  regressors, it shows the critical values from Table~IIc of Phillips and Ouliaris
  (1990) for $g - 1$ regressors.

\item This leaves the case where $\E(\Delta \mathbf{y}_{2t}) = \mathbf{0}$ and $\E(\Delta y_{1t}) \neq 0$. Since $y_{1t}$ has drift
  and $\mathbf{y}_{2t}$ does not, we need to include time in the regression (10.3.1) in order
  to remove a linear trend from the residuals. The discussion for the previous
  case makes it clear that the ADF $t$-statistic calculated using the residuals from
  a regression of $y_{1t}$ on a constant, $g$ driftless I(1) regressors $\mathbf{y}_{2t}$, \emph{and} time has
  the asymptotic distribution tabulated in Table~10.1(b) for $g + 1$ regressors (or
  Table~IIc of Phillips and Ouliaris, 1990, for $g$ regressors). For example, if
  $g = 2$, the regressor (10.3.1) has a constant, two I(1) regressors, and time; the
  critical values can be found from Table~10.1(b) for \emph{three} regressors.
\end{enumerate}

\begin{table}[htbp]
\centering
\caption{Critical Values for the ADF $t$-Statistic Applied to Residuals}
\label{tab:10.1}
\medskip
Estimated Regression: $y_{1t} = \mu + \boldsymbol{\gamma}'\mathbf{y}_{2t}$
\medskip

\begin{tabular}{l cccc}
\toprule
Number of regressors, & & & & \\
excluding constant & 1\% & 2.5\% & 5\% & 10\% \\
\midrule
\multicolumn{5}{l}{(a) Regressors have no drift} \\
\midrule
1 & $-3.96$ & $-3.64$ & $-3.37$ & $-3.07$ \\
2 & $-4.31$ & $-4.02$ & $-3.77$ & $-3.45$ \\
3 & $-4.73$ & $-4.37$ & $-4.11$ & $-3.83$ \\
4 & $-5.07$ & $-4.71$ & $-4.45$ & $-4.16$ \\
5 & $-5.28$ & $-4.98$ & $-4.71$ & $-4.43$ \\
\midrule
\multicolumn{5}{l}{(b) Some regressors have drift} \\
\midrule
1 & $-3.96$ & $-3.67$ & $-3.41$ & $-3.13$ \\
2 & $-4.36$ & $-4.07$ & $-3.80$ & $-3.52$ \\
3 & $-4.65$ & $-4.39$ & $-4.16$ & $-3.84$ \\
4 & $-5.04$ & $-4.77$ & $-4.49$ & $-4.20$ \\
5 & $-5.36$ & $-5.02$ & $-4.74$ & $-4.46$ \\
\bottomrule
\end{tabular}

\medskip
\small
\textsc{Source:} For panel~(a), Phillips and Ouliaris (1990, Table~IIb). For panel~(b), the
first row is from Fuller (1996, Table~10.A.2), and the other rows are from Phillips
and Ouliaris (1990, Table~IIc).
\end{table}

Several comments are in order about the residual-based test for cointegration.

\begin{itemize}
\item (Test consistency)\quad The alternative hypothesis is that $\mathbf{y}_t$ is cointegrated (i.e.,
  $h \geq 1$). The test is consistent against the alternative as long as $y_{1t}$ is cointegrated
  with $\mathbf{y}_{2t}$.\footnote{%
  For the case of $h = 1$, this is an implication of Theorem~5.1 of Phillips and Ouliaris (1990). Their remark~(d)
  to this theorem shows that the test is consistent when $h > 1$ and hence $\mathbf{y}_{2t}$ is cointegrated.}
  The reason (we will discuss it in more detail below for the case with
  $h = 1$) is that the OLS residuals from the regression (10.3.1) will converge to a
  stationary process. However, if $y_{1t}$ is I(1) and not part of the cointegration rela\-tionship,
  then the test may have no power against the alternative of cointegration
  because the OLS residuals, $y_{1t} - \boldsymbol{\gamma}'\mathbf{y}_{2t}$, with a nonzero coefficient (of unity) on
  the I(1) variable $y_{1t}$, will not converge to a stationary process. Thus, the choice
  of normalization (of which variable should be used as the dependent variable)
  matters for the consistency of the test.

\item (Should time be included in the regression?)\quad If time is included in the regres\-sion
  (10.3.1), then the drift in $y_{1t}$, $\E(\Delta y_{1t})$, affects only the time coefficient,
  making the numerical value of the residuals (and hence the ADF $t$-value) invari\-ant
  to $\E(\Delta y_{1t})$. This means that the case~3 procedure, which adds time to the
  regressors, can be used for case~1, where $\E(\Delta y_{1t})$ happens to be zero. That is,
  if you include time in the regression (10.3.1), then the appropriate critical value
  for case~1 is given from Table~10.1(b) for $g + 1$ regressors. The procedure is
  also valid for case~2, because if time is included in addition to a constant and
  $g$ I(1) regressors with drift, then the regression can be rewritten as a regression
  with a constant, $g$ driftless I(1) regressors, and time that combines the drifts in
  the $g$ I(1) regressors. This regression falls under case~3 and the critical values
  provided by Table~10.1(b) for $g + 1$ regressors apply. Therefore, when time is
  included in the regression (10.3.1), the same critical values can be used, regard\-less
  of the location of drifts. A possible disadvantage is reduced power in finite
  samples. The finite-sample power is indeed lower at least for the DGPs exam\-ined
  by Hansen (1992a) in his simulations.

\item (Choice of lag length)\quad The requirement (10.3.3) does not provide a practical
  rule in finite samples for selecting the lag length $p$ in the augmented autoregres\-sion
  to be estimated on the residuals. There seems to be no work in the context of
  the residual-based test comparable to that of Ng and Perron (1995) for univari\-ate
  I(1) processes. The usual practice is to proceed as in the univariate context,
  which is to use the Akaike information criterion (AIC) or the Bayesian infor\-mation
  criterion (BIC), also called the Schwartz information criterion (SIC), to
  determine the number of lagged changes.

\item (Finite-sample considerations)\quad Besides the residual-based ADF $t$-test, a num\-ber
  of tests are available for testing the null of no cointegration. They include
  tests proposed by Phillips and Ouliaris (1990), Stock and Watson (1988), and
  Johansen (1988). Haug (1996) reports a recent Monte Carlo study examining the
  finite-sample performance of these and other tests. It reveals a tradeoff between
  power and size distortions (that is, tests with the least size distortion tend to have
  low power). The residual-based ADF $t$-test with the lag length chosen by AIC,
  although less powerful than some other tests, has the least size distortion for the
  DGPs examined.
\end{itemize}

The following example applies the residual-based test with the ADF $t$-statistic
to the consumption-income relationship.

\begin{defbox}
\textbf{Example 10.5\quad (Are consumption and income cointegrated?):}\quad As was
already mentioned in connection with Figure~10.1(a), log income ($y_t$) and
log consumption ($c_t$) appear to be cointegrated with a cointegrating vector of
$(1, -1)$. This figure, however, is rather deceiving. The plot of $y_t - c_t$ (which
is the log of the saving rate) in Figure~10.1(b) shows an upward drift right
after the war and a downward drift since the mid 1980s. (The latter is the
well-publicized fact that the U.S.\ personal saving rate has been declining.)
As seen below, the test results depend on whether to include these periods
or not. We initially focus on the sample period of 1950:Q1 to 1986:Q4 (148
obervations).

\end{defbox}

\begin{figure}[htbp]
  \centering
  \includegraphics[width=0.9\textwidth,trim=50 350 50 350,clip]{figures/fig_10_1b.png}
  \caption{Log Income Minus Log Consumption}
  \label{fig:10.1b}
\end{figure}

\begin{defbox}
\noindent

We first test whether the two series are individually I(1), by conduct\-ing
the ADF $t$-test with a constant and time trend in the augmented auto\-regression.
In applying the BIC to select the number of lagged changes in
the augmented autoregression, we follow the same practice of the previous
chapter: the maximum length ($p_{max}$) is $\bigl[12 \cdot (T/100)^{1/4}\bigr]$ (the integer part of
$[12 \cdot (T/100)^{1/4}]$), the sample period is fixed at $t = p_{max} + 2, p_{max} + 3, \ldots, T$
in the process of choosing the lag length $p$, and, given $p$, the maximum sam\-ple
of $t = p + 2, p + 3, \ldots, T$ is used to estimate the augmented autoregres\-sion
with $p$ lagged changes. For the present sample size, $p_{max}$ is~13.

For disposable income, the BIC selects the lag length of~0 and the ADF
$t$-statistic ($t^\tau$) is $-1.80$. The 5 percent critical value from Table~9.2\,(c) is
$-3.41$, so we accept the hypothesis that $y_t$ is I(1). For consumption, the lag
length by the BIC is~1 and the ADF $t$ statistic is $-2.07$. So we accept the I(1)
null at 5 percent. Thus, both series might well be described as I(1) with drift.

Turning to the residual-based test for cointegration, the OLS estimate of
the static regression is
\begin{equation}
  c_t = -0.046 + 0.975\, y_t, \quad
  R^2 = 0.998, \quad
  t = \text{1950:Q1 to 1986:Q4}.
  \label{eq:10.3.4} \tag{10.3.4}
\end{equation}
\hfill{\small (0.009) \quad (0.0037)}

\medskip

To conduct the residual-based ADF $t$ test, an augmented autoregression with\-out
constant or time is estimated on the residuals with the lag length of~0
selected by the BIC. The ADF $t$ statistic is $-5.49$. Since the series have time
trends, we turn to Table~10.1(b), rather than Table~10.1(a), to find critical val\-ues.
For $g = 1$, the 5 percent critical value is $-3.41$, so we can reject the null
of no cointegration.

If the residual-based test is conducted on the entire sample of 47:Q1
to 98:Q1, the ADF $t$ statistic is $-2.94$ with the lag length of~1 determined
by BIC, and thus, we cannot reject the hypothesis that the consumption-on-income
regression is spurious!
\end{defbox}

\subsection*{Testing the Null of Cointegration}

In the above test, cointegration is taken as the alternative hypothesis rather than the
null. But very often in economics the hypothesis of economic interest is whether
the variables in question are cointegrated, so it would be desirable to develop tests
where the null hypothesis is that $h = 1$ rather than that $h = 0$. Very recently,
several tests of the null of cointegration have been proposed. For a catalogue of
such tests, see Maddala and Kim (1998, Section~4.5). As was true in the testing of
the null that a univariate process is I(0), there is no single test of cointegration used
by the majority of researchers yet.\footnote{%
The procedures of Johansen (1988) and Stock and Watson (1988) cannot be used for testing the null of
cointegration against the alternative of no cointegration, because in their tests the alternative hypothesis specifies
\emph{more} cointegration than the null.}

\subsection*{Question for Review}

\begin{enumerate}
\item For the residual-based test for cointegration with the regression (10.3.1) not
  including time as a regressor, verify that there is very little difference in the
  critical value between case~1 and case~2.
\end{enumerate}
